\section{Kernel mapping and observable imprints}
\label{app:kernels}

\subsection{From memory kernel to effective background}
Consider the background response
\begin{equation}
H^2(t) = \frac{8\pi G}{3}\rho(t) + \int_{0}^{t} K(\Delta t)\,H(t-\Delta t)\,d\Delta t,
\end{equation}
with a minimal resonant kernel
\begin{equation}
K(\Delta t)=\frac{\kappa}{\tau}\,e^{-\Delta t/\tau}\cos(\omega_0 \Delta t)\,,\qquad \Delta t\ge 0.
\end{equation}
Laplace transforming, $\mathcal{L}\{H\}(s)=\hat H(s)$, gives
\begin{equation}
\hat H(s)\;=\;\frac{\hat H_{\Lambda{\rm CDM}}(s)}{1-\hat K(s)}\,,\qquad
\hat K(s)=\frac{\kappa}{\tau}\,\frac{s+\tfrac{1}{\tau}}{(s+\tfrac{1}{\tau})^2+\omega_0^2}\,.
\end{equation}
For small $|\hat K|$, $H \simeq H_{\Lambda{\rm CDM}}\,(1+\hat K+\cdots)$, i.e.\ a controlled deformation of the background.

\subsection{Perturbations and $f\sigma_8$ echoes}
Linear matter growth obeys
\begin{equation}
\ddot\delta + 2H\dot\delta - 4\pi G\rho_m \delta \;=\; \int_{0}^{t}\! \tilde K(\Delta t)\,\delta(t-\Delta t)\,d\Delta t,
\end{equation}
where $\tilde K$ is an induced kernel. To leading order in small memory,
\begin{equation}
f\sigma_8(z)\;=\;[f\sigma_8]_{\Lambda{\rm CDM}}(z)\,\Big[1 + A\,e^{-{\Delta t}/{\tau}}\cos(\omega_0 \Delta t + \phi)\Big],
\end{equation}
with amplitude $A\propto \kappa$ and phase $\phi$ set by initial conditions. The redshift period is $\Delta z \simeq \pi/\omega_0$ to first order.

\subsection{Parameter roles}
\begin{itemize}
  \item $\tau$: sets decay envelope of the residual; larger $\tau$ $\Rightarrow$ longer-lived echoes.
  \item $\omega_0$: fixes the echo spacing $\Delta z \simeq \pi/\omega_0$.
  \item $\kappa$: overall amplitude sign/size; $\kappa\!\to\!0$ kills the signal.
  \item $Q=\omega_0\tau/2$: sharpness of the feature; controls phase coherence.
\end{itemize}
These partially orthogonal roles aid identifiability across RSD and lensing.
